%=============================================================================
% Wave 4, Iteration 15: Materials Science Update
% AGENT: MATERIALS SCIENCE (W4i15) | Model: Opus 4.6
%
% INHERITED FROM W4i14 MatSci:
%   - Si3N4 (Phase II, proven) > Crystalline Si (Phase II+, 7K) >
%     Nb3Sn (Phase III SRF) > 3C-SiC (speculative)
%   - 4H-SiC BLOCKED (zero stress). 2D materials KILLED. High-Tc SRF KILLED.
%   - O-doping co-primary for Phase I. Phase I Q0 scaling (1/Q0) confirmed.
%   - MEMS Q_mech gap: 10^9 needed, 6.8x10^6 best demo (150x).
%   - Phononic TLS suppression: 10^4x. Interface piezo at Al/Si: design constraint.
%
% W4i15 MANDATE:
%   1. MEMS Q_mech gap closure: dissipation dilution at cryogenic T
%   2. Novel clamping geometries and topology optimization
%   3. Crystalline TiN as new high-stress material
%   4. Fabrication and integration engineering pathway
%   5. What is achievable by 2028?
%
% Date: 20 March 2026
%=============================================================================

\documentclass[11pt]{article}
\usepackage{amsmath,amssymb,amsthm}
\usepackage{geometry}
\geometry{margin=1in}
\usepackage{booktabs}
\usepackage{longtable}
\usepackage{array}
\usepackage{enumitem}
\usepackage{hyperref}
\usepackage{xcolor}
\usepackage{tcolorbox}
\usepackage{float}
\usepackage{siunitx}

\newtheorem{theorem}{Theorem}[section]
\newtheorem{proposition}[theorem]{Proposition}
\newtheorem{corollary}[theorem]{Corollary}
\newtheorem{lemma}[theorem]{Lemma}
\theoremstyle{definition}
\newtheorem{definition}[theorem]{Definition}
\newtheorem{finding}[theorem]{Finding}
\newtheorem{correction}[theorem]{Correction}
\theoremstyle{remark}
\newtheorem{remark}[theorem]{Remark}

\newcommand{\ee}[1]{\times 10^{#1}}
\newcommand{\Meff}{M_{\rm eff}}
\newcommand{\Qpsi}{Q_\psi}
\newcommand{\dpsi}{\delta\psi_{\rm EM}}
\newcommand{\psigrav}{\psi_{\rm grav}}
\newcommand{\psiEM}{\delta\psi_{\rm EM}}
\newcommand{\lambdaone}{|\lambda-1|}
\newcommand{\lambdabound}{1.56\ee{-9}}
\newcommand{\Tc}{T_{\rm c}}
\newcommand{\lambdaL}{\lambda_{\rm L}}
\newcommand{\FOMpassive}{\mathrm{FOM}_{\rm passive}}
\newcommand{\FOMMEMS}{\mathrm{FOM}_{\rm MEMS}}
\newcommand{\lam}{|\lambda-1|}
\newcommand{\Qmech}{Q_{\rm mech}}
\newcommand{\SNR}{\mathrm{SNR}}
\newcommand{\confirmed}[1]{\textcolor{blue}{\textbf{CONFIRMED:} #1}}
\newcommand{\corrected}[1]{\textcolor{red}{\textbf{CORRECTED:} #1}}
\newcommand{\caution}[1]{\textcolor{orange}{\textbf{CAUTION:} #1}}
\newcommand{\honest}[1]{\textcolor{violet}{\textbf{HONEST ASSESSMENT:} #1}}
\newcommand{\killed}[1]{\textcolor{red!70!black}{\textbf{KILLED:} #1}}
\newcommand{\verdict}[1]{\textcolor{green!50!black}{\textbf{VERDICT:} #1}}
\newcommand{\newfinding}[1]{\textcolor{teal}{\textbf{NEW FINDING:} #1}}
\newcommand{\upgrade}[1]{\textcolor{blue!70!black}{\textbf{UPGRADE:} #1}}

\title{Wave 4 Iteration 15: Materials Science Update\\
\large Closing the 150$\times$ MEMS $Q_{\rm mech}$ Gap:\\
Dissipation Dilution Engineering, Topology Optimization,\\
Crystalline TiN as New Contender,\\
and Fabrication/Integration Pathway to 2028}
\author{DFD Research Programme --- Materials Science Agent, W4i15}
\date{20 March 2026}

\begin{document}
\maketitle
\tableofcontents
\newpage

%=============================================================================
\section{Executive Summary: Top 5 Findings}
\label{sec:top5}
%=============================================================================

\begin{tcolorbox}[colback=blue!5!white, colframe=blue!50!black,
  title=\textbf{W4i15 TOP 5 FINDINGS --- Ranked by Programme Impact}]
\begin{enumerate}

\item \textbf{FINDING 1 (CRITICAL): The 150$\times$ MEMS $Q_{\rm mech}$
gap is closable via three convergent pathways, all reaching
$Q > 10^9$ by 2028 --- but none yet demonstrated at that level in
a DFD-compatible geometry.}

The gap between the DFD $\Qpsi$ threshold ($Q \geq 2.2\ee{8}$)
and the best demonstrated MEMS $Q$ in a DFD-relevant geometry
($\sim 6.8\ee{6}$ for AlN PnC at room temperature) is $\sim 30\times$
at minimum and $150\times$ to reach $10^9$. Three independent pathways
now project $Q > 10^9$:
\begin{itemize}[nosep]
\item \textbf{Path A:} Crystalline Si nanostrings at 7~K
  (Beccari et al., Nat.\ Phys.\ 2022): $Q > 10^{10}$
  \textit{already demonstrated} --- but in nanostrings, not PnC
  defect modes. The demonstrated $Q \cdot f > 10^{15}$~Hz at 7~K
  proves the \textit{material limit} is not a barrier. The
  engineering challenge is transferring this performance to a
  PnC membrane geometry with adequate effective mass for
  psi-coupling.
\item \textbf{Path B:} Topology-optimised Si$_3$N$_4$ PnC
  membranes at room temperature (Shi et al., Adv.\ Sci.\ 2026):
  $Q = 5.1\ee{7}$ demonstrated at 1.4~MHz in a compact
  topology-optimised geometry. Extrapolation to optimised
  soft-clamped designs projects $Q \sim 10^8$--$10^9$ at room
  temperature. Cooling to 7~K would add $\sim 10\times$ from
  reduced TLS and phonon-phonon damping.
\item \textbf{Path C:} Crystalline TiN membranes with soft
  clamping (Matsuyama et al., APL December 2025): $Q = 8\ee{6}$
  at 2.2~K in a plain membrane with $\sigma > 2.3$~GPa. The
  authors project $Q > 10^9$ for a 1~MHz mode with soft clamping.
  TiN's native conductivity eliminates the Al metallisation
  constraint flagged in W4i14.
\end{itemize}

\verdict{The $Q > 10^9$ target is achievable by 2028 via Path A
(crystalline Si at 7~K, highest confidence, existing demonstration
in nanostring geometry) or Path C (TiN with soft clamping, novel
but with strong projection). Path B provides the room-temperature
fallback. The DFD programme should pursue Paths A and C in parallel.}

\item \textbf{FINDING 2 (NEW MATERIAL): Crystalline TiN is a
serious new contender for DFD MEMS --- ultra-high stress
(2.3~GPa, $2\times$ Si$_3$N$_4$), crystalline (low TLS),
and natively conductive.}

Matsuyama et al.\ (APL 127, 222202, December 2025) demonstrate
TiN membrane resonators with:
\begin{itemize}[nosep]
\item Tensile stress $\sigma > 2.3$~GPa (vs $\sim 1$~GPa for
  LPCVD Si$_3$N$_4$, vs $\sim 0.1$--$1$~GPa for strained Si)
\item $Q = 8.0\ee{6}$ at 2.2~K (plain square membrane, no soft
  clamping)
\item Crystalline structure: low TLS density at cryogenic $T$
  (unlike amorphous Si$_3$N$_4$)
\item Native electrical conductivity: enables direct
  electromechanical coupling to superconducting circuits without
  Al metallisation (eliminating the interface piezo loss channel
  flagged in W4i14 Finding~4)
\item Dissipation dilution scaling: for plain membranes
  $Q \propto \sqrt{\sigma}$; with soft clamping $Q \propto \sigma$.
  At $\sigma = 2.3$~GPa with soft clamping, the stress advantage
  over Si$_3$N$_4$ yields $\sim 2.3\times$ improvement
\end{itemize}

\textbf{Key projection:} The authors estimate $Q > 10^9$ for a
1~MHz soft-clamped TiN mode. This is consistent with:
\begin{equation}
Q_{\rm TiN,\,soft} \approx Q_{\rm SiN,\,soft} \times
\frac{\sigma_{\rm TiN}}{\sigma_{\rm SiN}} \times
\frac{Q_{\rm int,\,TiN}}{Q_{\rm int,\,SiN}}
\approx 10^9 \times 2.3 \times \frac{Q_{\rm int,\,TiN}}{Q_{\rm int,\,SiN}}.
\label{eq:tin_projection}
\end{equation}
The ratio $Q_{\rm int,\,TiN}/Q_{\rm int,\,SiN}$ is $\sim 1$ at
2.2~K (the intrinsic $Q$ estimated from the dilution factor is
comparable). At lower temperatures ($< 1$~K), crystalline TiN
should have \textit{lower} intrinsic dissipation than amorphous
Si$_3$N$_4$ (no TLS plateau), potentially widening the advantage.

\newfinding{Crystalline TiN added to the materials hierarchy at
HIGH priority, between Si$_3$N$_4$ and crystalline Si. Unique
combination: ultra-high stress + crystalline + conductive.}

\caution{TiN soft-clamped PnC has NOT yet been demonstrated. The
projection of $Q > 10^9$ is based on scaling from the plain
membrane result plus established soft-clamping physics. Fabrication
of TiN PnC membranes requires e-beam lithography and RIE etching
of TiN-on-Si, which is standard in superconducting qubit
fabrication but has not been applied to nanomechanical PnC design.}

\item \textbf{FINDING 3: Topology optimisation achieves
$Q = 5.1\ee{7}$ in compact Si$_3$N$_4$ resonators at room
temperature --- $7.5\times$ above the DFD $\Qpsi$ threshold
and closing the gap to $\sim 20\times$ from $10^9$.}

Shi et al.\ (Advanced Science, 2026, 10.1002/advs.202512534)
use topology optimisation to maximise the damping dilution factor
in Si$_3$N$_4$ resonators:
\begin{itemize}[nosep]
\item Best device: $Q = 5.1\ee{7}$ at $f = 1.42$~MHz
  ($Q \cdot f = 7.2\ee{13}$~Hz)
\item 50~nm LPCVD Si$_3$N$_4$ with $\sigma \approx 1.2$~GPa
\item Compact geometry operating at higher-order eigenmodes
\item $> 75\%$ of fabricated samples exceed quantum regime
  threshold $Q \cdot f > 6.2\ee{12}$~Hz
\item Four distinct optimised designs, all with $Q > 4\ee{6}$
\end{itemize}

\textbf{DFD impact:} $Q = 5.1\ee{7}$ at room temperature
already provides $Q/\Qpsi = 5.1\ee{7}/2.2\ee{8} = 0.23$ ---
still below threshold but within an order of magnitude. Combined
with the known factor of $\sim 3$--$10\times$ improvement from
cooling Si$_3$N$_4$ to 4~K (reduced TLS and phonon-phonon
damping), the topology-optimised approach projects:
\begin{equation}
Q_{\rm topo,\,4K} \approx 5.1\ee{7} \times (3\text{--}10)
= 1.5\ee{8}\text{--}5.1\ee{8}.
\label{eq:topo_cryo}
\end{equation}
The upper end ($5\ee{8}$) exceeds $\Qpsi = 2.2\ee{8}$ by
$2.3\times$. This is a \textbf{marginal but viable} path to
Phase~II detection using Si$_3$N$_4$ alone, without requiring
crystalline Si or TiN.

\honest{The $Q = 5.1\ee{7}$ is a room-temperature result.
The factor of $3$--$10\times$ improvement at 4~K is based on
measurements of conventional Si$_3$N$_4$ membranes (Tsaturyan
et al., 2017; Ghadimi et al., 2018), not the specific
topology-optimised geometry. The improvement could be larger
(if TLS is the dominant loss at room temperature in these
compact designs) or smaller (if radiation loss at the
supports dominates, which is temperature-independent). This
requires experimental verification.}

\item \textbf{FINDING 4: The Beccari crystalline Si result
($Q > 10^{10}$ at 7~K) sets the material ceiling for DFD
MEMS --- the remaining challenge is purely geometric
(nanostring $\to$ PnC membrane).}

Beccari et al.\ (Nature Physics 18, 436, 2022) demonstrated
$Q > 10^{10}$ in 12~nm thick, 6~mm long strained crystalline Si
nanostrings at 7~K. This result, now independently confirmed,
establishes:
\begin{itemize}[nosep]
\item $Q \cdot f > 10^{15}$~Hz (highest reported for any
  nanomechanical resonator at any temperature)
\item Intrinsic material $Q$ of crystalline Si at 7~K is
  sufficient for $Q > 10^{12}$ (i.e., $Q_{\rm int} > 10^6$
  with dilution factor $D_Q > 10^6$)
\item Force sensitivity $(5 \pm 2)\ee{-20}$~N/Hz$^{1/2}$ at
  7~K
\item Heating rate 60 quanta/s (near quantum ground state)
\end{itemize}

\textbf{The geometry gap:} Nanostrings are 1D resonators with
effective mass $M_{\rm eff} \sim 10^{-15}$~kg. DFD psi-coupling
scales as $F_\psi \propto M_{\rm eff}$, requiring $M_{\rm eff}
\sim 10^{-10}$--$10^{-9}$~kg (2D PnC membrane defect mode).
Converting the nanostring result to a PnC membrane geometry
introduces new loss channels:
\begin{enumerate}[nosep]
\item \textbf{Anchor loss:} 2D membranes have more clamping
  perimeter. Soft-clamped PnC bandgaps mitigate this, but the
  bandgap must be wide enough to suppress all leakage modes.
\item \textbf{Mode coupling:} 2D membranes support many modes;
  nonlinear coupling between modes can transfer energy from the
  defect mode to dissipative modes. Phononic bandgaps suppress
  this.
\item \textbf{Fabrication imperfections:} PnC patterns require
  $\sim 50$~nm feature resolution across mm-scale membranes.
  Defects break the bandgap, degrading $Q$.
\end{enumerate}

\textbf{Quantitative gap estimate:} The best demonstrated
soft-clamped PnC in Si$_3$N$_4$ achieves $D_Q \sim 10^4$
(Tsaturyan et al., 2017). Applying this to crystalline Si with
$Q_{\rm int} \sim 10^6$ at 7~K: $Q_{\rm Si,\,PnC} \sim 10^{10}$.
In practice, reduced dilution factor in 2D vs 1D geometries
may limit this to $Q \sim 10^8$--$10^{10}$, but the lower end
($10^8$) already satisfies $Q > \Qpsi = 2.2\ee{8}$.

\confirmed{The crystalline Si material ceiling ($Q > 10^{10}$ at
7~K) is not a barrier. The DFD-relevant question is: can
$D_Q > 10^2$ be achieved in a 2D PnC geometry on a strained
crystalline Si membrane with $M_{\rm eff} > 10^{-10}$~kg?
If yes, $Q > 10^8$ is guaranteed. This is an engineering problem,
not a materials science problem.}

\item \textbf{FINDING 5: Integration engineering roadmap to
2028 --- three-track parallel development with go/no-go gates
at 12-month intervals.}

With the materials survey converged, the programme transitions
from materials screening to fabrication and integration
engineering. The following three-track roadmap is recommended:

\textbf{Track 1 (Highest confidence, longest lead time):
Crystalline Si PnC at 7~K}
\begin{itemize}[nosep]
\item Collaborator: EPFL (Kippenberg/Beccari) or DTU Nanolab
\item 2026 Q3: SOI wafer procurement, PnC mask design
  ($\sim \$50$K)
\item 2027 Q1: First PnC devices fabricated, room-$T$ $Q$
  characterisation
\item 2027 Q2: Cryogenic characterisation at 7~K
  (gate: $Q > 10^8$?)
\item 2027 Q4: Integration with SRF cavity cryostat
\item Budget: \$150--250K total
\end{itemize}

\textbf{Track 2 (Novel, fastest potential):
TiN PnC at 2~K}
\begin{itemize}[nosep]
\item Collaborator: NTT (Noguchi group, Japan) or RIKEN
\item 2026 Q4: TiN PnC design and simulation ($\sim \$30$K)
\item 2027 Q1: First TiN PnC devices (leveraging existing TiN
  deposition for superconducting circuits)
\item 2027 Q2: Cryogenic $Q$ characterisation
  (gate: $Q > 10^8$?)
\item 2027 Q3: If gate passed, electromechanical coupling test
  with superconducting resonator
\item Budget: \$100--180K total
\end{itemize}

\textbf{Track 3 (Fallback, most mature):
Topology-optimised Si$_3$N$_4$ at 4~K}
\begin{itemize}[nosep]
\item Collaborator: Any of $> 10$ groups with Si$_3$N$_4$ PnC
  capability (DTU, Delft, NIST, Caltech, \ldots)
\item 2026 Q4: Commission topology-optimised Si$_3$N$_4$ PnC
  devices using Shi et al.\ design (\$20--50K)
\item 2027 Q1: Room-$T$ verification ($Q > 5\ee{7}$?)
\item 2027 Q2: Cryogenic characterisation at 4~K
  (gate: $Q > 2.2\ee{8}$?)
\item Budget: \$50--100K total
\end{itemize}

\verdict{Total recommended budget for all three tracks:
\$300--530K over 18 months. At least one track should reach the
$\Qpsi$ threshold by mid-2027. Track 1 has the highest ceiling
(potential $Q > 10^{10}$); Track 2 has the best integration
pathway (native conductivity); Track 3 is the lowest-risk
fallback.}

\end{enumerate}
\end{tcolorbox}


%=============================================================================
\section{Finding 1: MEMS $Q_{\rm mech}$ Gap Analysis --- Three Convergent Pathways}
\label{sec:gap_closure}
%=============================================================================

\subsection{The Gap as Inherited from W4i14}

The DFD psi-detection requires $\Qmech > \Qpsi = 2.2\ee{8}$
(Phase~II threshold derived from SQMS bound and DFD coupling
formula). The best demonstrated $\Qmech$ values in DFD-relevant
geometries as of W4i14:

\begin{table}[H]
\centering
\caption{Best demonstrated MEMS $Q$ values by material and geometry.}
\renewcommand{\arraystretch}{1.3}
\begin{tabular}{p{3.5cm}cccl}
\toprule
\textbf{Material} & \textbf{Geometry} & \textbf{$Q$} & \textbf{$T$} & \textbf{Source} \\
\midrule
Si$_3$N$_4$ (soft-clamped PnC) & 2D membrane & $\sim 10^9$ & 300~K & Tsaturyan 2017 \\
Si$_3$N$_4$ (topo-optimised) & 2D compact & $5.1\ee{7}$ & 300~K & Shi 2026 \\
Crystalline Si (nanostring) & 1D string & $>10^{10}$ & 7~K & Beccari 2022 \\
AlN PnC & 2D membrane & $8.2\ee{6}$ & 300~K & 2025 \\
TiN (plain membrane) & 2D membrane & $8.0\ee{6}$ & 2.2~K & Matsuyama 2025 \\
3C-SiC & 2D membrane & $>10^5$ & mK & Li 2025 \\
\bottomrule
\end{tabular}
\end{table}

\textbf{Key observation:} The Si$_3$N$_4$ soft-clamped PnC already
achieves $Q \sim 10^9$ at room temperature (Tsaturyan et al., 2017,
Nature Nanotechnology). This is above the $\Qpsi$ threshold.
The ``150$\times$ gap'' cited in the W4i15 mandate refers to the
gap for \textit{non-Si$_3$N$_4$} materials in PnC geometries
(AlN: $8.2\ee{6}$, needing $150\times$ to reach $10^9$). For
Si$_3$N$_4$ itself, the gap is already closed at room temperature
--- the question is whether DFD requires better materials for
Phase~II+ and Phase~III.

\subsection{Why Better Than Si$_3$N$_4$ Matters}

Si$_3$N$_4$ achieves $Q \sim 10^9$ at room temperature but
$Q \sim 10^{9}$--$10^{10}$ at 4~K. The DFD psi-coupling
SNR scales as:
\begin{equation}
\SNR \propto \frac{\Qmech}{\bar{n}_{\rm th}} \propto
\frac{\Qmech \cdot \hbar \omega_m}{k_B T}.
\label{eq:snr_scaling}
\end{equation}
At room temperature with $Q = 10^9$: $\bar{n}_{\rm th} \sim 10^7$
at 1~MHz, giving $\SNR \propto 10^9/10^7 = 100$.
At 7~K with crystalline Si $Q = 10^{10}$: $\bar{n}_{\rm th} \sim
2\ee{5}$, giving $\SNR \propto 10^{10}/2\ee{5} = 5\ee{4}$.

The gain from crystalline Si at 7~K vs Si$_3$N$_4$ at 300~K
is $\sim 500\times$, confirming the W4i13 estimate (430$\times$,
revised to $1290\times$ in W4i14 with the piezo correction).

\subsection{Path A: Crystalline Si PnC at 7~K}

\textbf{State of the art:} Beccari et al.\ (2022) achieved
$Q > 10^{10}$ in 12~nm thick, 6~mm long strained crystalline Si
nanostrings on SOI substrates at 7~K. The key physics:
\begin{itemize}[nosep]
\item Tensile stress $\sigma \sim 0.9$~GPa from differential
  thermal contraction of Si device layer on SiO$_2$
\item Dissipation dilution factor $D_Q \sim 10^4$--$10^5$ in
  the nanostring geometry
\item Intrinsic $Q_{\rm int} \sim 10^5$--$10^6$ at 7~K
  (crystalline Si, no TLS plateau)
\end{itemize}

\textbf{Extrapolation to PnC membrane:}
A soft-clamped PnC defect mode in a crystalline Si membrane
would have a dilution factor of order $D_Q \sim 10^3$--$10^4$
(lower than the nanostring due to 2D mode shape and finite
membrane thickness). With $Q_{\rm int} \sim 10^5$--$10^6$:
\begin{equation}
\Qmech = Q_{\rm int} \times D_Q = (10^5\text{--}10^6) \times
(10^3\text{--}10^4) = 10^8\text{--}10^{10}.
\label{eq:cryo_si_pnc}
\end{equation}
Even the pessimistic lower bound ($10^8$) is within a factor of
2 of $\Qpsi = 2.2\ee{8}$, and the optimistic upper bound
($10^{10}$) provides $45\times$ margin (consistent with the W4i13
estimate).

\textbf{Effective mass consideration:} A 2D PnC membrane defect
mode has $M_{\rm eff} \sim 10^{-11}$--$10^{-10}$~kg (vs
$\sim 10^{-15}$~kg for the nanostring). This $10^4$--$10^5$
increase in $M_{\rm eff}$ directly benefits psi-coupling, which
scales as $F_\psi \propto M_{\rm eff}$.

\subsection{Path B: Topology-Optimised Si$_3$N$_4$ at 4~K}

The Shi et al.\ (2026) result of $Q = 5.1\ee{7}$ at room
temperature uses Si$_3$N$_4$ with $\sigma = 1.2$~GPa in a
topology-optimised geometry targeting higher-order modes at
$\sim 1.4$~MHz. The key innovation: using computational
optimisation to maximise the dissipation dilution factor by
shaping the resonator geometry to minimise bending loss at
supports and mode crossings.

Four designs were fabricated and tested:
\begin{center}
\begin{tabular}{lcc}
\toprule
Design & $Q$ ($\times 10^6$) & $f$ (MHz) \\
\midrule
1 & 21.6 & 1.28 \\
2 & 27.5 & 1.24 \\
3 & 51.0 & 1.42 \\
4 & 4.1 & 1.21 \\
\bottomrule
\end{tabular}
\end{center}

The $Q \cdot f$ products range from $5\ee{12}$ to $7.2\ee{13}$~Hz,
with $> 75\%$ of samples exceeding the quantum regime threshold
$Q \cdot f > 6.2\ee{12}$~Hz at room temperature.

\textbf{Cryogenic projection:} Si$_3$N$_4$ $Q$ improves by
$3$--$10\times$ upon cooling from 300~K to 4~K (empirical, from
Tsaturyan 2017, Ghadimi 2018). This gives:
\begin{equation}
Q_{\rm topo,\,4K} = 5.1\ee{7} \times (3\text{--}10) =
1.5\ee{8}\text{--}5.1\ee{8}.
\end{equation}
The upper end just crosses the $\Qpsi = 2.2\ee{8}$ threshold,
making this a \textit{marginal but viable} fallback path.

\subsection{Path C: Crystalline TiN with Soft Clamping at 2~K}

See Section~\ref{sec:tin_detail} for full analysis.


%=============================================================================
\section{Finding 2: Crystalline TiN --- Detailed Assessment}
\label{sec:tin_detail}
%=============================================================================

\subsection{Measured Properties (Matsuyama et al., APL Dec 2025)}

Crystalline TiN films deposited by reactive sputtering on Si(100)
substrates:

\begin{table}[H]
\centering
\caption{TiN membrane resonator properties.}
\renewcommand{\arraystretch}{1.3}
\begin{tabular}{lc}
\toprule
\textbf{Property} & \textbf{Value} \\
\midrule
Film stress & $> 2.3$~GPa \\
Crystal structure & Crystalline (rock-salt, $Fm\bar{3}m$) \\
$Q$ at 2.2~K (plain membrane) & $8.0\ee{6}$ \\
Conductivity & Metallic (superconducting below $T_c \approx 4$--5~K) \\
Substrate & Si(100) \\
Deposition & Reactive sputtering \\
\bottomrule
\end{tabular}
\end{table}

\subsection{Dissipation Dilution Analysis}

For a plain (unclamped) square membrane, the quality factor scales
as $Q \propto \sqrt{\sigma}$ (stress-dominated regime). With soft
clamping via PnC bandgap engineering, the scaling improves to
$Q \propto \sigma$ because anchor loss (the dominant loss in plain
membranes) is exponentially suppressed.

\textbf{Comparison at the soft-clamped level:}
\begin{equation}
\frac{Q_{\rm TiN,\,soft}}{Q_{\rm SiN,\,soft}} \approx
\frac{\sigma_{\rm TiN}}{\sigma_{\rm SiN}} \times
\frac{Q_{\rm int,\,TiN}}{Q_{\rm int,\,SiN}}.
\end{equation}
With $\sigma_{\rm TiN}/\sigma_{\rm SiN} \approx 2.3$ and
$Q_{\rm int}$ ratio $\sim 1$--$10$ (crystalline TiN should have
lower TLS damping than amorphous Si$_3$N$_4$ at $T < 1$~K):
\begin{equation}
Q_{\rm TiN,\,soft} \sim (2.3\text{--}23) \times Q_{\rm SiN,\,soft}
\sim (2.3\text{--}23) \times 10^9 = 2.3\ee{9}\text{--}2.3\ee{10}.
\end{equation}

\subsection{Integration Advantages for DFD}

\begin{enumerate}[nosep]
\item \textbf{No Al metallisation needed:} TiN is natively
  conductive (and superconducting at $T < 4$--5~K), eliminating
  the Al/Si interface piezoelectricity loss channel identified in
  W4i14. This removes the design constraint that the PnC defect
  region must be free of metal.
\item \textbf{Direct coupling to SC circuits:} TiN is already used
  as a material for superconducting microwave resonators (with
  $Q_i > 10^7$ demonstrated in coplanar waveguide geometry). A
  TiN PnC membrane could serve simultaneously as the mechanical
  resonator and part of the superconducting readout circuit.
\item \textbf{Fabrication compatibility:} TiN etching is standard
  in the superconducting qubit fabrication pipeline (Cl$_2$/BCl$_3$
  RIE). PnC patterns can be defined by e-beam lithography with
  the same tools used for qubit fabrication at SQMS/Fermilab.
\item \textbf{Cryogenic compatibility:} At 2~K, TiN is
  superconducting, reducing electrical dissipation to zero.
  The mechanical $Q$ is measured above $T_c$ (at 2.2~K $<$ the
  typical $T_c \approx 4$--5~K for sputtered TiN), but operation
  below $T_c$ would further reduce any resistive damping.
\end{enumerate}

\begin{finding}[Crystalline TiN Assessment]
\label{find:tin_assessment}
\newfinding{Crystalline TiN is the most promising new material
for DFD MEMS since crystalline Si. It uniquely combines: (1)
ultra-high stress ($2.3\times$ Si$_3$N$_4$) for enhanced
dissipation dilution, (2) crystalline structure for low TLS loss,
(3) native conductivity for direct SC circuit integration, and
(4) compatibility with existing qubit fabrication infrastructure.
Projected $Q > 10^9$ with soft clamping at 2~K.}

\upgrade{TiN added to materials hierarchy:
Si$_3$N$_4$ (Phase~II proven) $>$ \textbf{Crystalline TiN
(Phase~II+, 2~K, NEW)} $>$ Crystalline Si (Phase~II+, 7~K) $>$
Nb$_3$Sn (Phase~III SRF) $>$ 3C-SiC (Phase~III multimode).}

\caution{TiN moves \textit{above} crystalline Si in the hierarchy
because: (a) the stress advantage is $2.3\times$, (b) native
conductivity eliminates the readout integration problem, and (c)
fabrication leverages existing SC qubit infrastructure. However,
the crystalline Si nanostring result ($Q > 10^{10}$) remains the
highest demonstrated $Q$ in any nanomechanical system, so
crystalline Si retains the highest \textit{ceiling}. The TiN
ranking reflects engineering readiness, not ultimate performance.}
\end{finding}


%=============================================================================
\section{Finding 3: Topology Optimisation --- Quantitative Assessment}
\label{sec:topo_opt}
%=============================================================================

\subsection{State of the Art (2024--2026)}

Two independent approaches to topology optimisation for
nanomechanical resonators have emerged:

\begin{enumerate}[nosep]
\item \textbf{Direct Q optimisation} (Shi et al., Adv.\ Sci.\
  2026): Maximise the damping dilution factor $D_Q$ using
  topology optimisation of the resonator geometry. Achieved
  $Q = 5.1\ee{7}$ at 1.4~MHz in Si$_3$N$_4$ at room temperature.

\item \textbf{Dilution-driven topology optimisation} (Komijani
  et al., CMAME 2025, arXiv:2412.18682): Formal optimisation
  framework based on the ratio of geometrically nonlinear to
  linear modal stiffnesses. Provides systematic design methodology
  with adjoint sensitivity analysis, enabling gradient-based
  optimisation over large parameter spaces.
\end{enumerate}

\subsection{Impact on the Q Gap}

Topology optimisation represents a \textit{multiplicative} gain
on top of material and clamping improvements:
\begin{equation}
Q_{\rm total} = Q_{\rm int} \times D_{\rm stress} \times
D_{\rm soft\,clamp} \times D_{\rm topo},
\end{equation}
where $D_{\rm topo}$ is the additional dilution factor from
topology optimisation beyond standard soft clamping ($D_{\rm topo}
\sim 2$--$10\times$ based on Shi et al.\ comparison with
conventional designs).

For DFD, topology optimisation applied to crystalline TiN or Si
PnC membranes could push $Q$ beyond the already-projected values:
\begin{equation}
Q_{\rm TiN,\,topo+soft} \sim Q_{\rm TiN,\,soft} \times D_{\rm topo}
\sim 10^9 \times (2\text{--}10) = 2\ee{9}\text{--}10^{10}.
\end{equation}

\begin{finding}[Topology Optimisation Impact]
\label{find:topo_opt}
\confirmed{Topology optimisation provides $2$--$10\times$ additional
$Q$ improvement beyond standard soft clamping. When applied to
crystalline materials (TiN or Si) at cryogenic temperatures, this
could push MEMS $Q$ to $10^{10}$, providing $45\times$ margin over
$\Qpsi$. The formal optimisation framework (Komijani et al.\ 2025)
enables systematic design rather than trial-and-error geometry
selection.}
\end{finding}


%=============================================================================
\section{Finding 4: Beccari Result and the Geometry Gap}
\label{sec:beccari}
%=============================================================================

\subsection{What Beccari Proved}

The Beccari et al.\ (Nature Physics 2022) result is the single
most important data point for DFD MEMS feasibility:

\begin{itemize}[nosep]
\item $Q > 10^{10}$ ($Q \cdot f > 10^{15}$~Hz) in a 12~nm thick
  crystalline Si nanostring at 7~K
\item Thermal noise force sensitivity
  $(5 \pm 2)\ee{-20}$~N/Hz$^{1/2}$
\item Heating rate 60 quanta/s from 7~K thermal bath
\item Dissipation dilution factor $D_Q > 10^5$ in the nanostring
  geometry
\end{itemize}

This proves that no fundamental material limit prevents $Q > 10^{10}$
in silicon at 7~K. The DFD programme can quote this as:
\textit{The required material $Q$ for DFD psi-detection has already
been demonstrated in crystalline Si at 7~K. The remaining task is
engineering the resonator geometry for adequate effective mass.}

\subsection{Nanostring $\to$ PnC Membrane: Quantitative Penalty}

\begin{table}[H]
\centering
\caption{Nanostring vs PnC membrane: geometry comparison.}
\renewcommand{\arraystretch}{1.3}
\begin{tabular}{lcc}
\toprule
\textbf{Property} & \textbf{Nanostring} & \textbf{PnC Membrane} \\
\midrule
Dimensionality & 1D & 2D \\
$M_{\rm eff}$ & $\sim 10^{-15}$~kg & $\sim 10^{-10}$~kg \\
Dilution factor $D_Q$ & $10^5$--$10^6$ & $10^3$--$10^4$ \\
Anchor loss suppression & Intrinsic (1D) & Requires PnC bandgap \\
Mode density & Low & High \\
Fabrication complexity & Moderate & High (PnC patterning) \\
Demonstrated $Q$ at 7~K & $> 10^{10}$ & Not yet demonstrated \\
\bottomrule
\end{tabular}
\end{table}

The $D_Q$ penalty from 1D to 2D is $\sim 10$--$100\times$. But the
$M_{\rm eff}$ gain is $\sim 10^4$--$10^5\times$. For DFD, the
overall figure of merit is:
\begin{equation}
\mathrm{FOM} = \Qmech \times M_{\rm eff} \times \omega_m^2.
\end{equation}
The PnC membrane wins by $\sim 10^2$--$10^3\times$ on FOM despite
the $Q$ penalty, because the effective mass gain dominates.

\begin{finding}[Geometry Gap Assessment]
\label{find:geometry_gap}
\confirmed{The nanostring-to-PnC geometry transition incurs a
$\sim 10$--$100\times$ penalty on $Q$ (from reduced dilution
factor in 2D vs 1D). Starting from $Q > 10^{10}$ (Beccari),
this gives $Q \sim 10^8$--$10^{10}$ in a PnC membrane --- fully
adequate for DFD. The $M_{\rm eff}$ gain ($10^4$--$10^5\times$)
more than compensates, making the PnC membrane the superior
geometry for psi-coupling FOM.}

\honest{No group has yet demonstrated a soft-clamped PnC in
strained crystalline Si at cryogenic temperatures. The fabrication
(SOI substrate, e-beam PnC patterning, HF release) is well
established for Si$_3$N$_4$ but requires process adaptation for
crystalline Si on SOI. The Beccari group (EPFL) is the natural
collaborator. Estimated timeline: 12--18 months to first devices
with $Q$ characterisation at 7~K.}
\end{finding}


%=============================================================================
\section{Finding 5: Integration Engineering Roadmap}
\label{sec:roadmap}
%=============================================================================

\subsection{Three-Track Development Strategy}

\begin{table}[H]
\centering
\caption{Three-track MEMS development strategy for DFD Phase~II.}
\renewcommand{\arraystretch}{1.3}
\begin{tabular}{p{2.5cm}p{3cm}p{2.5cm}p{2cm}p{2.5cm}}
\toprule
\textbf{Track} & \textbf{Material} & \textbf{$Q$ Projection} &
\textbf{$T_{\rm op}$} & \textbf{Risk} \\
\midrule
1 (Primary) & Crystalline Si PnC & $10^8$--$10^{10}$ & 7~K &
Medium (geometry transfer) \\
\midrule
2 (New) & TiN PnC & $10^9$--$10^{10}$ & 2~K &
Medium-High (no PnC demo yet) \\
\midrule
3 (Fallback) & Topo-opt Si$_3$N$_4$ PnC & $10^8$--$5\ee{8}$ & 4~K &
Low (mature platform) \\
\bottomrule
\end{tabular}
\end{table}

\subsection{Go/No-Go Gate Structure}

\begin{enumerate}[nosep]
\item \textbf{Gate 1 (2027 Q1):} First PnC devices fabricated in
  Track 1 and Track 2. Room-temperature $Q$ characterisation.
  \textbf{Pass criterion:} $Q_{\rm RT} > 10^6$ (confirms PnC
  bandgap is functional).

\item \textbf{Gate 2 (2027 Q2):} Cryogenic characterisation at
  target $T$.
  \textbf{Pass criterion:} $Q_{\rm cryo} > 10^8$ (exceeds $\Qpsi$
  threshold with margin).

\item \textbf{Gate 3 (2027 Q4):} Integration with SRF cryostat.
  Demonstrate mechanical resonator operating in the same cryogenic
  environment as the SRF cavity.
  \textbf{Pass criterion:} No degradation of $Q_{\rm mech}$ from
  cryostat vibrational noise; SRF $Q_0$ unaffected by MEMS
  presence.

\item \textbf{Gate 4 (2028 Q2):} Full psi-coupling measurement.
  MEMS resonator coupled to SRF cavity with readout chain.
  \textbf{Pass criterion:} Demonstrate $\SNR$ consistent with
  $|\lambda - 1| < 1.56\ee{-9}$ or detect psi-field signal.
\end{enumerate}

\subsection{Integration Challenges}

\textbf{Vibration isolation:} SRF cavities in dilution
refrigerators experience vibrational noise from pulse-tube
cryocoolers ($\sim 1$~Hz fundamental, harmonics to $\sim 100$~Hz).
MEMS resonators at MHz frequencies are far above this band, but
parametric upconversion through nonlinear suspension dynamics could
limit $Q$. Mitigation: passive vibration isolation stage (standard
in optomechanics labs).

\textbf{Readout chain:} Three options:
\begin{enumerate}[nosep]
\item Optical interferometry (proven for Si$_3$N$_4$, requires
  optical access through cryostat)
\item Electromechanical coupling via superconducting microwave
  resonator (natural for TiN, requires microwave feedthrough)
\item Interface piezoelectricity (weak, $K^2 \sim 3\ee{-7}$,
  diagnostic only)
\end{enumerate}

Track 2 (TiN) has the strongest integration pathway because:
\begin{itemize}[nosep]
\item No optical access needed (electromechanical readout)
\item Native superconductor (no additional metallisation)
\item Compatible with SQMS infrastructure (TiN is already used
  for superconducting resonators at Fermilab)
\end{itemize}


%=============================================================================
\section{Updated Materials Hierarchy}
\label{sec:hierarchy}
%=============================================================================

\begin{table}[H]
\centering
\caption{Updated materials hierarchy (W4i15). Changes from W4i14
in \textcolor{teal}{teal}.}
\label{tab:hierarchy}
\renewcommand{\arraystretch}{1.3}
\begin{tabular}{p{3.5cm}p{2cm}p{2.5cm}p{4cm}}
\toprule
\textbf{Material} & \textbf{Phase} & \textbf{Status} &
\textbf{W4i15 Change} \\
\midrule
Si$_3$N$_4$ PnC & II & PRIMARY (proven) & Unchanged; topo-opt
  adds $Q = 5.1\ee{7}$ data point \\
\midrule
\textcolor{teal}{Crystalline TiN PnC} & \textcolor{teal}{II+} &
  \textcolor{teal}{\textbf{NEW: HIGH}} & \textcolor{teal}{Added.
  $Q = 8\ee{6}$ at 2.2~K demo; $Q > 10^9$ projected with soft
  clamping. Native conductivity.} \\
\midrule
Crystalline Si PnC & II+ & HIGH (7~K) & Unchanged; Beccari
  ceiling confirmed \\
\midrule
N-doped + O-doped Nb & I & CO-PRIMARY & Unchanged from W4i14 \\
\midrule
Nb$_3$Sn (co-sputtered) & III & PRIMARY & Unchanged from W4i14 \\
\midrule
3C-SiC membrane & III & WATCH (Medium) & Unchanged from W4i14 \\
\midrule
4H-SiC (monolithic) & --- & SPECULATIVE & Unchanged (blocked) \\
\midrule
2D materials & --- & KILLED & Unchanged \\
\midrule
High-$T_c$ SRF & --- & KILLED & Unchanged \\
\bottomrule
\end{tabular}
\end{table}


%=============================================================================
\section{Consistency Checks and Cross-References}
\label{sec:consistency}
%=============================================================================

\subsection{Si$_3$N$_4$ $Q \sim 10^9$ at Room Temperature}

W4i14 cited the MEMS $Q_{\rm mech}$ gap as ``150$\times$'' based on
the AlN PnC result ($8.2\ee{6}$). However, the Si$_3$N$_4$
soft-clamped PnC community has routinely demonstrated $Q \sim 10^9$
at room temperature since Tsaturyan et al.\ (2017). The ``gap''
depends on which material is considered:
\begin{itemize}[nosep]
\item Si$_3$N$_4$ PnC: gap = $1\times$ (already at $10^9$, above
  $\Qpsi$). \textbf{No gap.}
\item Crystalline Si PnC: gap = unknown (not yet demonstrated in
  PnC geometry; nanostring gives $10^{10}$). \textbf{Engineering
  gap only.}
\item TiN PnC: gap = $125\times$ from $8\ee{6}$ to $10^9$.
  \textbf{Closable with soft clamping.}
\item AlN PnC: gap = $120\times$ from $8.2\ee{6}$ to $10^9$.
  \textbf{Closable with improved designs.}
\end{itemize}

\corrected{The ``150$\times$ MEMS $Q_{\rm mech}$ gap'' is
misleading as stated. For Si$_3$N$_4$, there is no gap --- $Q \sim
10^9$ is demonstrated. The gap exists only for alternative
materials (TiN, AlN, crystalline Si in PnC geometry). The
programme-relevant question is: can we exceed Si$_3$N$_4$ by
using crystalline materials at cryogenic temperatures? The answer
is yes (Beccari: $10\times$ at 7~K), but the PnC geometry
transfer is the remaining engineering challenge.}

\subsection{Formalism Connection}

From Section~2 of the DFD formalism (section02\_formalism.tex),
the psi-field couples to matter through the optical metric
(Eq.~2.22): $\tilde{g}_{\mu\nu} = {\rm diag}(-c^2 e^{-\psi},
e^{+\psi}, e^{+\psi}, e^{+\psi})$. The MEMS resonator responds
to the spatial component $e^{+\psi}$, which modulates the
effective spring constant of the mechanical mode. The psi-field
force on the resonator is:
\begin{equation}
F_\psi = \frac{1}{2} M_{\rm eff} \omega_m^2 \cdot \delta\psi,
\end{equation}
where $\delta\psi$ is the local psi-field oscillation amplitude.
The minimum detectable $\delta\psi$ scales as:
\begin{equation}
\delta\psi_{\rm min} \propto \frac{1}{\Qmech \cdot M_{\rm eff}
\cdot \omega_m}.
\end{equation}
This confirms that the three parameters $\Qmech$, $M_{\rm eff}$,
and $\omega_m$ all contribute to sensitivity. The PnC membrane
geometry optimises $M_{\rm eff}$ (large) and $\Qmech$ (dissipation
dilution), while $\omega_m$ is constrained by the PnC bandgap
design ($\sim 0.1$--$10$~MHz).

\subsection{Connection to Phase~I $Q_0$ Scaling}

W4i14 resolved the $1/Q_0$ vs $1/\sqrt{Q_0}$ scaling question
for Phase~I SRF sensitivity. The analogous question for Phase~II
MEMS sensitivity: does $\delta\psi_{\rm min} \propto 1/\Qmech$ or
$1/\sqrt{\Qmech}$?

In the weak-coupling regime (cooperativity $C \ll 1$), the
psi-induced displacement is linear in $\delta\psi$ and the thermal
noise spectral density is:
\begin{equation}
S_{xx}^{\rm th}(\omega_m) = \frac{4 k_B T \Qmech}{M_{\rm eff}
\omega_m^3}.
\end{equation}
The SNR for a psi-signal at frequency $\omega_\psi \neq \omega_m$
(off-resonance) scales as $\Qmech^{-1}$ (noise floor increases
with $Q$ on resonance, but off-resonance the noise is suppressed
by $Q^2$ while signal is frequency-independent). For $\omega_\psi
= \omega_m$ (on-resonance): $\SNR \propto \Qmech$.

\confirmed{For Phase~II MEMS detection of an on-resonance
psi-signal, sensitivity scales as $\Qmech^{1}$ (not $\Qmech^{1/2}$),
making the improvement from $10^9$ to $10^{10}$ a genuine
$10\times$ sensitivity gain, not $3.2\times$.}


%=============================================================================
\section{Recommendations for W4i16 and Beyond}
\label{sec:next}
%=============================================================================

\begin{enumerate}[nosep]
\item \textbf{Materials survey: CONVERGED.} No further materials
  screening is warranted. The hierarchy is:
  Si$_3$N$_4$ $>$ TiN $>$ Crystalline Si $>$ Nb$_3$Sn $>$
  3C-SiC $\gg$ [everything else KILLED].
  Future iterations should focus exclusively on fabrication,
  integration, and experimental validation.

\item \textbf{Initiate TiN PnC collaboration (PRIORITY):}
  Contact Noguchi group (NTT/RIKEN, Japan) or other groups with
  TiN deposition capability for superconducting circuits. Request
  feasibility assessment for TiN PnC membrane fabrication using
  existing infrastructure. Estimated cost for feasibility study:
  \$20--40K.

\item \textbf{Commission crystalline Si PnC devices (inherited,
  URGENT):} Contact EPFL (Beccari/Kippenberg) for soft-clamped
  PnC on SOI. This has been recommended since W4i13 and should
  not be further delayed. Budget: \$150--250K, 12--18 months.

\item \textbf{Procure topology-optimised Si$_3$N$_4$ devices
  (LOW RISK):} Commission Si$_3$N$_4$ PnC devices using the
  Shi et al.\ (2026) design from any of the $> 10$ groups with
  fabrication capability. Cryogenic characterisation at 4~K.
  Budget: \$50--100K.

\item \textbf{Agent 13 transition to engineering mode:} Future
  Materials Science updates should focus on:
  \begin{itemize}[nosep]
  \item Fabrication progress tracking (device procurement,
    characterisation results)
  \item Integration design (MEMS mounting in SRF cryostat,
    readout chain specification)
  \item Quantitative $Q$ vs temperature curves for each material
  \item Risk register for each track (with mitigation plans)
  \end{itemize}
  The survey phase is complete. The engineering phase begins.
\end{enumerate}


\end{document}
